\documentclass[ecta,nameyear,draft]{econsocart}
\usepackage{amsmath,amsfonts,amssymb,amsthm,booktabs,array,longtable}
\usepackage{graphicx}
\usepackage[T1]{fontenc}
\usepackage{mathpazo,microtype,needspace,xurl}
\usepackage{xr-hyper}
\RequirePackage[colorlinks,citecolor=blue,linkcolor=blue,urlcolor=blue]{hyperref}
\startlocaldefs
\newtheorem{theorem}{Theorem}
\newtheorem{proposition}[theorem]{Proposition}
\newtheorem{lemma}[theorem]{Lemma}
\newtheorem{corollary}[theorem]{Corollary}
\newcommand{\E}{\mathbb E}
\newcommand{\R}{\mathbb R}
\endlocaldefs
\setlength{\emergencystretch}{3em}
\makeatletter\def\copyright@text{Prepared for referee review}\makeatother

\setcounter{secnumdepth}{2}

% Stable current-document identity for clean typesetting reproduction.
\ifdefined\pdfinfoomitdate\pdfinfoomitdate=1\fi
\ifdefined\pdfsuppressptexinfo\pdfsuppressptexinfo=15\fi
\ifdefined\pdftrailerid\pdftrailerid{}\fi

\externaldocument[][nocite]{build/supp}[supp.pdf]
\externaldocument[][nocite]{build/appendix}[appendix.pdf]
\begin{document}
\begin{frontmatter}
\title{Neural Bellman Operators}
\runtitle{Neural Bellman Operators}
\begin{aug}
\author[id=au1,addressref={add1}]{\fnms{Qian}~\snm{QI}\ead[label=e92ECTA]{qiqian@pku.edu.cn}}
\address[id=add1]{Peking University}
\end{aug}
\begin{abstract}
Neural Bellman Operators combine continuation critics with feasible actors for
recursive economic decisions. We characterize the complete action information
in paid scalar Bellman supports through an exact finite convex compatibility
problem, constructive interpolation and rational residual certificates. A
separate precision theorem connects action localization, original-policy regret,
value width and complete oracle work. A prospectively frozen experiment in the
original continuous-risk investment economy retains 56 neural and conventional
actors, strong online controls and dynamic programming. Adaptive precision
certifies 84 additional policy gains and meets the stronger regret target in
231 of 232 services. It uses 79.3 percent fewer extra integration contexts than
fixed precision, but takes 12.4 percent more complete time. Independent rational
evaluation confirms every accepted gain. Exact policy identities sharpen
conditional adoption comparisons. Matched hinge controls preclude attributing
shared geometry to neural architecture. The complete-support theorem and
executed acquisition service have distinct algorithmic scopes.
\end{abstract}
\begin{keyword}
\kwd{Dynamic programming}\kwd{Neural approximation}\kwd{Recursive risk}
\kwd{Policy qualification}\kwd{Computational economics}
\end{keyword}
\end{frontmatter}
\section{Introduction}\label{sec:introduction125}

An investment decision depends on the continuation values of the actions
that remain feasible, on the risk of their future consequences, and on the
cost of obtaining those values. A small average prediction error need not
identify a good investment. An accurate certificate for a fixed policy does
not by itself improve that policy. Neural Bellman Operators address this
economic decision problem by constructing a continuation critic together
with a feasible actor for a specified recursive objective. Their numerical
content must be assessed at the level of that complete decision service.

Our application is a constrained investment economy with a nonlinear
transition, a state-dependent capacity, a continuous innovation, and a nested
tail-risk preference. The primitive Bellman objective, rather than a fitted
surrogate or a discretized alternative economy, defines optimality. We ask
three related questions. What information about the next action is already
contained in a paid Bellman calculation? When should a service obtain a new
action observation or a more accurate evaluation of an existing one? Can a
neural continuation provide enough economic value to pay for its acquisition,
training, qualification and deployment relative to strong conventional
methods with the same information?

The first contribution is an exact finite description of the original-action
information in a complete scalar Bellman support record. Existing
qualification calculations return a global quadratic minorant as well as
an interval for the value at the queried action. Compressing this record to
a value and a derivative interval discards restrictions on other actions.
We retain every quadratic minorant. After removing the known strong
convexity, each minorant becomes an affine obstacle. Adjacent smooth-convex
interpolation conditions, together with a bounded-domain conjugate
inequality at each node and obstacle, are necessary and sufficient for
compatibility with all the original records. The proof constructs a
continuous piecewise-affine derivative realizing the information. A finite
convex program therefore gives the exact worst compatible action contrast
under this explicitly stated information interface. A rational residual
correction validates a bound even when an optimization routine supplies an
inexact candidate. This is a completion of the information used by the
Bellman certificate, with no claim that convex interpolation or dual
residual certification is itself new.

The second contribution links action localization to the complete cost of
continuous-risk information. The inherited midpoint rule already bounds
bracket contraction. We make explicit the support precision required for a
declared contraction, the sum of all refinement costs, the constrained
action-rounding allowance, and the separate terminal requirements for
original-policy regret and absolute objective width. Under a stated
tunable-oracle cost assumption, the resulting precision schedule has a
finite work bound; the assertion does not presume that a fixed experimental
precision ladder reaches the required level. When the available information
cannot certify a change, the algorithm preserves the last qualified feasible
policy and charges the work that was attempted.

These developments retain the original neural construction. Its critic is
built backwards from its own continuation and its actor carries a feasible
action witness. The operator foundations establish finite-horizon value
existence and uniqueness, measurable decisions, stability and implementation
allowances. One-sided and centered-risk accounts connect local errors to the
original recursive objective. The supplement supplies complete proofs and
the conditional bridge to stochastic control and nonlinear recursive
utility. The Brownian bridge requires its own well-posed backward equation,
regularity and stopping-boundary conditions. A result under expectation,
nested AVaR, or a value-dependent recursive driver retains that particular
objective and information filtration.

The numerical study asks whether the new acquisition service improves the
decision after paying for a strong original-model control. It uses all
56 previously acquired actors, covering neural ReLU, exactly matched
ordinary hinge, quadratic, spline, polynomial and grid representations,
and includes direct online and uniform-action dynamic-programming controls.
All economic comparisons retain the same continuous innovation law,
original future controller and service tolerances. Four states and the
complete method catalogue were fixed in a remotely recorded source
boundary before the first new evaluation. This deterministic study measures
action information and complete work on that specified census. Earlier
independent random-state evidence and its simultaneous inference remain
separately identified. Repeated process clocks are descriptive measurements,
not additional economic observations.

The strongest historical comparisons constrain the interpretation of the
new results. Exact ReLU and matched hinge encodings return the same
decisions. Previous complete-cost neural adoption domains were empty, and
traditional quadratic methods often won. We preserve these results and
their failed or unresolved computations. The present support theorem
applies equally to conventional and neural proposals; its rational strict
witness also has a cheap direct-support comparator. It demonstrates that
discarded information can affect a certified economic gain, not a unique
neural advantage. The original-economy experiment implements the separately
specified acquisition service; it is not an executed performance test of
the new complete-support convex program.

The contribution sought by Neural Bellman Operators remains a competitive
method for recursive economic decisions. The current results identify
precisely which parts of that contribution are established: feasible
own-future construction, original-objective certificates, exact action
information, and an accountable mechanism for acquiring further precision.
Whether a learned neural representation can deliver a positive complete
economic advantage over equally informed conventional representations
remains an empirical and theoretical question, addressed by the comparisons
and explicit adoption conditions rather than assumed from the architecture.

The article first specifies the operator and original economy. It then
develops the paid-support and complete-information results, the precision
and work account, and the current economic experiment. The technical
supplement contains the complete proof chains. A separate development
companion preserves the earlier models, statements, executions, negative
results and corrections, with source identities and a content-retention map.

\section{The economic operator and its policy account}\label{sec:core49}
Consider dates $t=0,\ldots,T$ and an invariant compact state domain $X\subset\R^d$. At date $t$, the action belongs to a nonempty compact set $A_t(x)$. The transition is $F_t(x,a,Z_{t+1})$, the current cost is $c_t(x,a)$, and the terminal cost is $g(x)$. Write $B_0=1$ and $B_t=\prod_{j<t}\beta_j$, with $\beta_j>0$. An implemented policy includes its observation rule, action selection and feasible numerical realization, not just its stored parameters.

For a continuation $v$, define
\begin{equation}\label{eq:operator49}
 Q_t(v;x,a)=c_t(x,a)+\beta_t\rho_t(v(F_t(x,a,Z))),\qquad
 \mathcal T_tv(x)=\inf_{a\in A_t(x)}Q_t(v;x,a).
\end{equation}
Throughout the deterministic comparison, the evaluations $\rho_t$ are monotone and cash invariant on a common domain containing all functions and constant translates used below. They are therefore nonexpansive in the supremum norm. Continuous objectives, effective compact feasible correspondences and measurable minimizers, or measurable arbitrarily near-minimizers with their error charged, are assumed. The stipulated conditional recursion gives Bellman comparison against all admissible adapted controls. In the earlier expected-cost experiments $\rho_t$ is conditional expectation, costs are bounded and innovations have their original continuous uniform law. Section~\ref{sec:continuous95} evaluates the same investment technology under an explicitly nested AVaR preference and retains the original continuous innovation law. The direct-cost inference in Section~\ref{sec:decisions49} is specific to expectation; it is not automatically an inference theorem for nonlinear recursive utility or games.

A Neural Bellman Operator is a construction and evaluation procedure for \eqref{eq:operator49}. It works backwards with its own acquired continuation, returns a neural critic and a feasible actor, and records an economic error allowance. Its continuation may be fitted or enclosed by an explicitly paid primitive Bellman calculation; an exact optimal continuation is never supplied as a free oracle. The constructive theorem below uses its own fitted $f_{t+1}$ at every date. In the new two-date application, the terminal function $g$ is an original primitive, and the paid last-date operator $T_1g$ supplies outward target intervals shared with every conventional comparator. The deployed last-date policy retains its separately proved error. Arbitrary trained networks can be evaluated by the first result below. The constructive backend specifies how a network and its actor are obtained without assuming successful nonconvex optimization.

\subsection{Action-dependent errors and complete-policy loss}
Let $H(a)$ be a true current cost and $\widehat H(a)$ its continuation approximation. If $\widehat a$ is $\delta$-optimal for $\widehat H$ on the same feasible action set, then
\begin{equation}\label{eq:center49}
 H(\widehat a)-\inf_aH(a)\le\delta+
 \operatorname{osc}_{a}\{\widehat H(a)-H(a)\}.
\end{equation}
Indeed, compare $\widehat a$ with an arbitrarily near-optimal true action and cancel the approximate objectives. The infimum limit gives the result. Additive continuation levels cancel; action-dependent discrepancies do not. Under a change in future policies, the continuation must be transported to the new economic problem before this inequality licenses reuse. The full development edition retains the original quantitative transport, recursive-domain and economic-application results. No transport bound is presumed valid after an unaccounted change of constraints or preferences.

\begin{proposition}[One-sided policy sandwich]\label{prop:sandwich49}
Suppose, uniformly in state,
\[
 f_t-\mathcal T_tf_{t+1}\le u_t,\qquad
 \mathcal T_t^\pi f_{t+1}-f_t\le v_t,
 \qquad \ell_T\le f_T-g\le u_T.
\]
Then the actual feasible policy satisfies
\begin{equation}\label{eq:sandwich49}
 0\le J_0^\pi-V_0^*\le
 \sum_{t<T}B_t(u_t+v_t)+B_T(u_T-\ell_T).
\end{equation}
The bound is invariant to date-specific constant shifts of the fitted continuations, provided the associated residual accounts are shifted consistently.
\end{proposition}
\noindent The proof is in Supplement, Section~\ref{supp:proof-core49-1}.


A sharper valid residual account for the same continuation and the same deployed actor sharpens knowledge of the policy. It does not retrain that policy or change its realized cost. Conversely, subtracting two upper bounds in \eqref{eq:sandwich49} does not identify the difference in actual costs. These distinctions govern both the historical comparisons and the new experiments.

\section{Constructing and compiling a feasible neural policy}\label{sec:construction49}
Assume the primitive objective has, for feasible state--action pairs, state modulus $A_t^Q=C_{x,t}+\beta_tF_{x,t}L_{t+1}$ and action modulus $D_t=C_{a,t}+\beta_tF_{a,t}L_{t+1}$. Let
\[
 d_H(A_t(x),A_t(z))\le K_t^A\|x-z\|_1,
 \qquad L_t=A_t^Q+D_tK_t^A,\qquad L_T=L_g.
\]
The Hausdorff condition and the primitive moduli imply that $\mathcal T_tf_{t+1}$ is $L_t$-Lipschitz whenever $f_{t+1}$ is $L_{t+1}$-Lipschitz. All moduli and numerical integration procedures are supplied by the economic primitives, not fitted from successful outcomes.

At state nodes $x_{t,i}$ with covering radius $h_t$, form effective feasible $k_t$-nets. Evaluate their own-future objectives with certified error at most $e_t$. Minimize the computed labels and retain the corresponding \emph{original} action $a_{t,i}$. A Borel repair $\Gamma_{t,i}(x)$ is feasible at $x$ and satisfies
\begin{equation}\label{eq:repair49}
 \|\Gamma_{t,i}(x)-a_{t,i}\|_1\le K_t^A\|x-x_{t,i}\|_1+\zeta_t.
\end{equation}
The allowance $\zeta_t$ concerns displacement; feasibility itself may not fail. Define
\begin{equation}\label{eq:cone49}
 f_t(x)=\min_i\{y_{t,i}+L_t\|x-x_{t,i}\|_1\},\quad
 i_t(x)=\min\arg\min_i\{y_{t,i}+L_t\|x-x_{t,i}\|_1\}.
\end{equation}
Return $\pi_t(x)=\Gamma_{t,i_t(x)}(x)$. Terminal labels approximate $g$ with error $e_T$ and use the same cone construction with modulus $L_g$.

\begin{theorem}[Own-future construction with feasible witnesses]\label{thm:construct49}
The functions in \eqref{eq:cone49} have exact affine--ReLU realizations and Lipschitz moduli $L_t$ independent of label error. The returned policy is feasible and measurable. It satisfies
\begin{align}\label{eq:enclosures49}
 -e_t\le f_t-\mathcal T_tf_{t+1}&\le D_tk_t+e_t+2L_th_t,\\
 \mathcal T_t^\pi f_{t+1}-f_t&\le e_t+D_t\zeta_t,\notag
\end{align}
and hence
\begin{equation}\label{eq:constructionloss49}
 J_0^\pi-V_0^*\le
 \sum_{t<T}B_t(D_tk_t+2e_t+2L_th_t+D_t\zeta_t)
 +B_T(2e_T+2L_gh_T).
\end{equation}
With effective covers, repairs and certified queries at the prescribed accuracy, allocating each nonnegative term a share of a positive economic tolerance gives a terminating finite construction. A fixed numerical precision or a capped finite experimental ladder need not contain that allocation.
\end{theorem}
\noindent The proof is in Supplement, Section~\ref{supp:proof-core49-2}.


The Lipschitz-extension principle is classical \citep{mcshane1934}. A native min-plus envelope and its exact ReLU circuit, using the same labels, witnesses and ties, are \emph{one method with two encodings}. They return the same function and policy. Neither Theorem~\ref{thm:construct49} nor a faster representation proves a uniquely neural capability or a training-success probability for Adam or L-BFGS. The mathematical increment is the complete feasible, own-future construction and implementation account; the empirical question is its work and policy value relative to conventional generators.


\section{Reference-centered risk kernels and complete-policy qualification}\label{sec:centeredrisk92}
A datewise worst-case error account can charge incompatible states at successive dates. Under nonlinear recursive preferences, replacing that account by an expectation under the physical transition law is not valid. We instead retain the conditional risk evaluation and center its dual representation at the actual stored NBO continuation. The resulting state-dependent certificate has an original-optimum interpretation and a finite occupation representation. It extends, rather than replaces, the signed qualification of Section~\ref{sec:risk91}.

\subsection{The centered operator}
Retain the original augmented economic state, feasible actions, nonnegative discounts, and bounded complete-policy recursions. At each $(t,x,a)$ suppose the original conditional evaluation admits the representation
\begin{equation}\label{eq:dualrisk92}
 \rho_{t,x,a}(v)=\sup_{q\in\mathcal Q_{t,x,a}}
 \left\{\int v(y)q(dy)-\psi_{t,x,a}(q)\right\}.
\end{equation}
Every $q$ is a probability measure on the original successor state space. The representation is assumed on a common bounded-function domain containing all continuations below, with the measurability and selection conditions needed for the original Bellman recursions. This is a hypothesis about the original preference primitive, not permission to replace a static risk objective by a nested one. Conditional risk duality and risk-averse dynamic programming are established principles \citep{ruszczynski2010risk,ruszczynski2014erratum}. For a nonconvex or non-cash-invariant aggregator lacking \eqref{eq:dualrisk92}, the earlier order and modulus results remain applicable.

Fix one qualified parameter $\theta$ and suppress its subscript. Its entire sequence of stored continuations is $f=(f_t)_{t=0}^T$, and its supplied complete feasible policy is $\pi$. Define, using the same $f_{t+1}$ in every term,
\begin{align}\label{eq:centered-def92}
 r_t(x,a)&=c_t(x,a)+\beta_t\rho_{t,x,a}(f_{t+1})-f_t(x),\notag\\
 k_t^f(x,a,q)&=\psi_{t,x,a}(q)+\rho_{t,x,a}(f_{t+1})
                         -\int f_{t+1}\,dq,\notag\\
 \mathcal E_{t,x,a}^f(z)&=\sup_{q\in\mathcal Q_{t,x,a}}
                   \left\{\int z\,dq-k_t^f(x,a,q)\right\}.
\end{align}
The penalty is nonnegative and has infimum zero, whether or not the supremum in \eqref{eq:dualrisk92} is attained. Direct substitution gives
\begin{equation}\label{eq:exact-center92}
 \mathcal E_{t,x,a}^f(z)
 =\rho_{t,x,a}(f_{t+1}+z)-\rho_{t,x,a}(f_{t+1}).
\end{equation}
In particular, $\mathcal E^f(0)=0$, and $\mathcal E^f$ is monotone, cash invariant, and nonexpansive in the uniform norm. Discarding $k^f$ is an upper relaxation, not this identity. Freezing one risk-maximizing kernel at $f$ is in general not even an upper relaxation.

\begin{theorem}[State-dependent original-optimum certificate]\label{thm:centeredrisk92}
Suppose verified functions satisfy $r_t^-(x,a)\le r_t(x,a)$ for every original feasible action and $r_t(x,\pi_t(x))\le r_t^{\pi,+}(x)$. Let bounded measurable $L_t,U_t$ satisfy
\begin{align}\label{eq:potential-cert92}
 L_T&\le g-f_T\le U_T,\notag\\
 L_t(x)&\le\inf_{a\in A_t(x)}
       \{r_t^-(x,a)+\beta_t\mathcal E_{t,x,a}^f(L_{t+1})\},\notag\\
 U_t(x)&\ge r_t^{\pi,+}(x)
                +\beta_t\mathcal E_{t,x,\pi_t(x)}^f(U_{t+1}).
\end{align}
Then, at every initial state,
\begin{equation}\label{eq:potential-loss92}
 f_0+L_0\le V_0^*\le V_0^\pi\le f_0+U_0,
 \qquad 0\le V_0^\pi-V_0^*\le U_0-L_0.
\end{equation}
The same statements hold at every intervening date and state. They are uniform over a qualified parameter class whenever all premises and the reported bound hold for each member with its own common continuation sequence and complete policy. The certificate permits endogenous states and actions changing at every date; it does not require a finite state space or independent deployment contexts.

With exact residuals and terminal discrepancies, equality recursions in \eqref{eq:potential-cert92} give $L_t=V_t^*-f_t$ and $U_t=V_t^\pi-f_t$. With scalar bounds from \eqref{eq:signed-input91}, their corresponding constant lower and upper sequences are feasible certificates here. Thus the exact state-dependent recursion is no looser than the signed scalar account. A strictly smaller bound can result without changing a fit, a feasible action, or an economic observation.
\end{theorem}

The last assertion concerns the best exact recursion or a verified improvement upon the constant certificates. It is not a claim that every approximately computed $L,U$ improves the scalar account. In implementation, the minimum of valid policy supersolutions and the maximum of valid optimum subsolutions retain validity, so an inferior new bound need not displace the existing one.

The lower inequality requires every original action. A discrete action search without a verified original-menu lower bound cannot establish it. For a fixed action, a feasible risk kernel supplies a lower bound on $\mathcal E^f(L)$; an upper bound requires control of the entire risk envelope. Numerical integration, feasible repairs, residual errors and penalty errors belong in these inequalities. If penalty enclosures are used, lower penalty endpoints enter an upper certificate and upper endpoints enter a lower witness. Approximate dual optimizers are accepted only through these independently checked inequalities.


\section{Continuous-law qualification in the investment model}\label{sec:continuous95}
The risk theorem becomes a numerical method only when its inequalities can be constructed from economic primitives. We now supply this construction for the original investment technology with a nested tail-risk preference. States and innovations remain continuous, every feasible action is covered, and the implemented terminal controller rather than a hypothetical exact optimizer determines policy cost. This application joins continuation construction, the original-optimum comparison and a representation comparison on one problem. It does not replace the general NBO construction or the other economic applications.

\subsection{Primitives and the object being approximated}
Let $d\geq2$ be even, $X=[0,1]^d$, and interpret subscripts cyclically. At dates $t=0,1$, the feasible investment is $0\leq a\leq A(x)=1/8+\bar x/8$, where $\bar x=d^{-1}\sum_jx_j$. Write $b_j=1/2$ for odd $j$ and $b_j=1/4$ for even $j$, and $s_j=(-1)^{j+1}$. The original transition is
\begin{equation}\label{eq:investment95}
 F_j(x,a,z)=\frac1{16}+\frac{x_j}{2}+\frac{x_{j+1}}8
       +\frac{x_j(1-x_{j+1})}{16}+b_ja+s_jz,
 \qquad z\sim\operatorname{Unif}[-1/32,1/32].
\end{equation}
Define
\begin{align}\label{eq:cost95}
 S(x)&=\frac2d\sum_j(x_j-5/8)^2+
       \frac1{4d}\sum_j(x_j-x_{j+1})^2+2(1/2-2\bar x)_+^2,\\
 g(x)&=S(x)+\frac2d\sum_j(x_j-5/8)^2,\qquad h(a)=a^2+4a^4.\nonumber
\end{align}
We use $\beta=15/16$ and, separately, $\alpha\in\{1/2,3/4\}$. Average value at risk is the upper-tail cost functional
\[
 \rho_\alpha(v)=\sup\left\{\int vq\,dP:
             0\leq q\leq(1-\alpha)^{-1},\ \int q\,dP=1\right\}.
\]
Here $P$ is the original uniform innovation law. The objective is recursively defined by $V_2=g$ and
\[
 V_t(x)=S(x)+\inf_{0\leq a\leq A(x)}
       \{h(a)+\beta\rho_\alpha(V_{t+1}(F(x,a,\cdot)))\}.
\]
It is not the static tail risk of total realized cost. A complete policy value replaces each infimum by the action prescribed by that policy, with the same recursive risk primitive.

\paragraph{Assumptions for the continuous construction.}
For this subsection, the primitives are exactly \eqref{eq:investment95}--\eqref{eq:cost95}; $\alpha$ is fixed before construction; the fitted continuation $f$ is a finite convex function on $X$; and every numerical lower or upper value is an outward enclosure. Feasible tail densities and feasible actions are retained explicitly. Terminal selection is the measurable, dyadically acquired controller described below, with its fallback included in the error bound. These concrete conditions imply compact feasible graphs, bounded costs, measurable minimizers, and the conditional-risk dual used in Theorem~\ref{thm:centeredrisk92}. They do not assert that every recursive preference or every trained network has this convex structure.

\begin{proposition}[Primitive curvature and state supports]\label{prop:curvature95}
For either risk level and every even $d$, the one-date objective
$Q_1(x,a)=S(x)+h(a)+\beta\rho_\alpha(g(F(x,a,\cdot)))$
is jointly convex on $X\times[0,1/4]$, with curvature bounded below by
\begin{equation}\label{eq:curvature95}
 \frac{4-5\beta/4}{d}\|v\|_2^2+2w^2
\end{equation}
in a direction $(v,w)$. Its value $V_1$ is convex on $X$ and satisfies
\begin{equation}\label{eq:lipschitz95}
 |V_1(x)-V_1(y)|\leq\frac{15}{d}\|x-y\|_1.
\end{equation}
Let $W$ be the objective evaluated under any fixed feasible risk density, at an anchor $(x_0,a_0)$, and let $(p_x,p_a)$ be its gradient there. For any $\lambda,\nu\geq0$, set $e=p_a+\lambda-\nu$. Then the following is a global affine minorant of $V_1$:
\begin{align}\label{eq:affine95}
 \ell(y)={}&W+\lambda(a_0-A(x_0))-\nu a_0-\frac{e^2}{4}\nonumber\\
 &+\left(p_x-\frac{\lambda}{8d}\boldsymbol{1}\right)^{\!\top}(y-x_0).
\end{align}
Neither optimal multipliers nor an exact optimizing action are required.
\end{proposition}

The curvature is obtained from the original nonlinear technology, not imposed on an interpolated value grid. The negative second-order term from $x_j(1-x_{j+1})$ is bounded using $|\partial_jg|\leq10/d$; the stage curvature dominates it. The allowance $e^2/4$ pays for imperfect stationarity. The proof in the Technical Supplement also establishes the global state modulus in \eqref{eq:lipschitz95}, including the state-dependent action constraint.

\subsection{An all-action certificate under the original innovation law}
Partition the innovation interval into $n$ equal bins $B_i$, with midpoints $z_i$ and probabilities $p_i=1/n$. An innovation bin is a set over which an inequality holds, not a replacement point mass in the economic model. Choose nonnegative $q_i\leq p_i/(1-\alpha)$ summing to one. These masses define a feasible density constant on each bin. For a fixed initial state and feasible proposal $a_0$, let $W_0$ be a lower bound on the objective under this density, obtained from globally valid affine supports of the continuation at the successor midpoints. Let $p$ enclose the derivative of the same supporting objective at $a_0$. Then
\begin{equation}\label{eq:actionlower95}
 \underline V_0(x)=\inf_{0\leq u\leq A(x)}
         \{W_0+p(u-a_0)+(u-a_0)^2\}
\end{equation}
is a lower bound on the original Bellman optimum when the supported continuation is $V_1$. The scalar quadratic is minimized over the full original interval. Interval uncertainty in $p$ is covered by taking the smaller minimum at its two endpoints, since the minimum is concave in $p$.

For a convex fitted continuation, midpoint supports give the lower risk witness and the larger endpoint value bounds the continuation throughout each bin. For $V_1$, Proposition~\ref{prop:curvature95} supplies midpoint supports and the allowance $15/(32n)$ covers each half-bin. Upper tail risk is the average of the largest $n(1-\alpha)$ upper bin bounds. This average is the exact risk of a dominating step function, hence an upper bound under the continuous law.

\begin{theorem}[Constructed original-optimum qualification]\label{thm:continuous95}
Suppose the supports and bin bounds above are outwardly verified. Suppose further that the complete terminal controller satisfies
$V_1\leq V_1^{\widehat\pi}\leq V_1+\varepsilon_1$ everywhere on $X$.
For every queried initial state, the lower value \eqref{eq:actionlower95} and the binwise upper bound $\overline J_0^{\widehat\pi}$ satisfy
\begin{equation}\label{eq:originalgap95}
 0\leq J_0^{\widehat\pi}(x)-V_0(x)
       \leq\overline J_0^{\widehat\pi}(x)-\underline V_0(x).
\end{equation}
For any stored convex continuation $f$, the same lower and upper constructions are admissible bounds for the centered-risk comparison of Theorem~\ref{thm:centeredrisk92}. Their validity does not depend on the representation of $f$, on a successful fit, or on the proposal being an exact minimizer.
\end{theorem}

To see the connection, write a verified continuation bracket as $\ell\leq V_1\leq u$ and use errors $L=\ell-f$, $U=u-f$. The centered identity gives
\[
 \rho_\alpha(f)+\mathcal E^f(L)=\rho_\alpha(\ell),\qquad
 \rho_\alpha(f)+\mathcal E^f(U)=\rho_\alpha(u).
\]
Joint evaluation avoids subtracting and then adding two separately rounded copies of $\rho_\alpha(f)$. The bound \eqref{eq:actionlower95} uses one feasible risk density to minorize the required all-action infimum; the policy upper bound controls the full risk envelope. Thus the application constructs sufficient inequalities for the original theorem. It does not claim that a non-neural continuation would receive a different validity theorem, or that computing the original-value witness is free.

The terminal controller observes each state coordinate rounded down to $40$ binary places. It accepts its floating proposal only after an outward all-action gap of at most $2^{-21}$ at that observed state; otherwise it uses rational bisection. The original-state transfer is less than $30\,2^{-40}$. Consequently $\varepsilon_1=2^{-20}$ covers the complete controller, including observation and fallback. Policy evaluation encloses this controller at every possible successor, not only at the midpoint states used by the integrator.


\paragraph{Well-posed original decisions.}
The comparison uses a specified conditional nested recursion on the
original observation filtration.  A sufficient primitive model has a
compact state space, nonempty compact feasible sets varying continuously
with state, continuous costs and invariant transitions, compact
innovations with their original fixed laws, and finite monotone
cash-invariant risk evaluations.  These assumptions make the Bellman
operator preserve continuous functions.  For implemented Borel policies,
the risk evaluation must also preserve Borel parameter kernels; the
supplement proves this property for expectation and AVaR.
Proposition~\ref{prop:finitefoundations119} proves, rather than
assumes, the resulting finite-horizon value existence and uniqueness,
attained Borel minimizers, comparison against all feasible adapted
history policies, and the dated stability bounds.  These are classical
foundations \citep{ruszczynski2010risk,ruszczynski2014erratum},
not a claim of a new fixed-point principle.  Pointwise transition
continuity with an unbounded innovation support is insufficient for
arbitrary monetary risk; the exact essential-uniform alternative is
stated in the proof.

The original investment primitives satisfy these conditions directly:
$F(x,a,z)\in[1/32,27/32]^d$, $1/8\le A(x)\le1/4$,
and all costs are bounded.  Their joint curvature gives
$V_1$ state strong-convexity modulus $(4-5\beta/4)/d$ and
root action modulus two.  Both dates therefore have unique continuous
constrained optimal actors.  Proposition~\ref{prop:investmentfoundations119}
also derives the continuous value gradient with its capacity multiplier,
the original-risk root derivative, and the actual observation/feasibility
allowance.  Differentiating a rounded actor is not used to establish
any of these claims.  The supplement records the full smooth critic--actor
graph and distinguishes its gradient from a frozen-target update.

The separate Brownian recursive-utility bridge has its own existence
and verification assumptions.  Proposition~\ref{prop:stoppedbsde119}
supplies the stopped-domain version, charging the original stopping
boundary and terminal condition.  Its Lipschitz driver can depend on
the backward value, in which case cash invariance is not automatic.
The continuous-time stability factor and smooth critic graph are
therefore verified within that original recursive criterion.

\subsection{Original marginal regularity and the acquired neural critic}
\label{sec:regularity122}

The present computation specializes the preceding investment primitives to
$d=2$. Write $q_x(a)=S(x)+h(a)+\beta\rho_\alpha
\{V_1(F(x,a,Z))\}$. Theorem~\ref{thm:rootmarginal121} proves that
$q_x$ has a continuous inward derivative on its full feasible interval,
is $2$-strongly convex, and has derivative Lipschitz constant
\begin{align*}
 L_\alpha&=5+\beta\left\{\frac{5M_\alpha}{16}
          +\frac{(4M_\alpha/5)^2\alpha}{(181/64)(1-\alpha)}\right\},\\
 M_\alpha&=\frac{1413}{128}
   +\beta\left\{\frac{127}{16}
                 +\frac{19625\alpha}{20480(1-\alpha)}\right\}.
\end{align*}
These constants are derived from the original technology and recursive risk,
including binding capacity and an adverse tail touching either innovation
boundary. They are not estimates of a fitted critic's curvature.

The retained acquired actor uses an anchored continuation of the form
\[
 \widehat R_\theta(x,v)=\widehat R_0(x)
       +\sum_j\theta_j\{\phi_j(x,v)-\phi_j(x,0)\},\qquad
 v=a/A(x).
\]
Its known stage cost remains analytic. Paid original-risk value enclosures
supply either the value loss or outward finite-difference marginal labels.
The complete finite-box fitting procedure and exact optimization-gap audits
are in Supplement, Section~\ref{sec:marginal121}. Their off-design
feature-cover theorem states the additional information needed to infer
global action accuracy from such training; the earlier experiment did not
execute that global cover. The new service below instead acquires and reuses
original information at the action actually proposed by the actor.

The resulting complete-policy cost remains
\[
 J_x(a)=q_x(a)+\eta_x(a),\qquad
 0\le\eta_x(a)\le\delta=\beta 2^{-20}.
\]
The allowance is uniform and concerns the original feasible terminal
controller. No derivative or convexity of $\eta_x$ is assumed. A
state-dependent value anchor must still be carried when a critic is exported
to an earlier Bellman operator: Theorem~\ref{thm:recursiveanchor121} gives
the full level-and-contrast recursion and its distinct noncash-invariant
version. Local action invariance never licenses deleting that recursive level.

% New R122 section. All inputs concern the original Q, not a fitted critic.
\section{Turning a paid Bellman certificate into a safe decision}
\label{sec:supportaction122}

A learned actor can be useful even when its fitted marginal is not a reliable
local direction. The original qualification calculation already purchases a
lower supporting objective and an upper original-value enclosure. The missing
step is to determine exactly which original action derivatives that information
permits, before paying for a second action or a global feature cover. This
section gives that conversion and a local policy safeguard. The calculation
is available without alteration to every conventional representation in the
comparison menu. Its contribution is an action-level use of already paid
Bellman information, not a neural-exclusive property of smooth descent.
Approximate supporting information and its accuracy--work tradeoff have a
classical inexact-oracle interpretation \citep{DGN2014Inexact122}.
Sharp smooth-convex information bounds also have an established interpolation
literature \citep{THG2017Interpolation122}. The calculation here states exactly
what the saved original-risk support licenses at a capacity-constrained action;
it does not claim a new general theory of projected gradient methods.

Fix an original state $x$, put $A=A(x)$, and write $q=Q_x$ on $[0,A]$.
The retained original-model result gives $\mu=2$ and
$L=L_\alpha$ in Theorem~\ref{thm:rootmarginal121}; no fitted-critic
curvature is used. Suppose a completed original certificate at $a\in[0,A]$
provides numbers $W,U,p_-,p_+$ and an actual supporting slope
$p\in[p_-,p_+]$ for which
\begin{equation}\label{eq:paidinformation122}
 q(a)\le U,\qquad
 q(b)\ge W+p(b-a)+\frac\mu2(b-a)^2\quad(0\le b\le A).
\end{equation}
The slope need not equal $q'(a)$. In particular, the slope of an approximate
terminal Bellman minorant is not differentiated as if its optimization error
were zero. Put $e=U-W\ge0$ and $c=(L-\mu)/2$.
All these data refer to the original optimal-future objective. A saved upper
bound on implemented-policy cost also bounds $q(a)$ and may be used
conservatively. Subtracting the future-policy allowance is permitted only
when the source formula establishes that the allowance was explicitly
added to an otherwise valid $q$ upper bound.

For $r>0$ define the computable uncertainty function
\begin{equation}\label{eq:psi122}
 \Psi(e,c,r)=\inf_{0<h\le r}\{e/h+ch\}
 =\begin{cases}
  2\sqrt{ec},&c>0,\ e\le cr^2,\\
  e/r+cr,&e>cr^2\text{ or }c=0.
 \end{cases}
\end{equation}
The first expression has value zero at $e=0$. When $r=0$, set
$\Psi(e,c,0)=+\infty$; absence of a feasible direction does not provide
an absent derivative bound. A rational or outward implementation can use
any feasible $h$ in place of the infimum.

\begin{theorem}[Exact derivative information in a paid quadratic support]
\label{thm:supportgradient122}
Suppose $q$ is continuously differentiable, $\mu$-strongly convex, and has
$L$-Lipschitz derivative, and \eqref{eq:paidinformation122} holds.
Then
\begin{equation}\label{eq:gradientinterval122}
 p_- -\Psi(e,c,A-a)\ \le q'(a)\le
 p_+ +\Psi(e,c,a).
\end{equation}
Intersecting with any separately proved original-gradient enclosure remains
valid. For the investment problem $[-7,7]$ is such an enclosure.
For an exact supplied slope $p$ and without that additional gradient
information, both endpoints of \eqref{eq:gradientinterval122} are sharp
in the stated information class: each finite endpoint is attained by a
function with constant curvature $L$, the same support, and $q(a)=U$.
At a missing side, the corresponding derivative is unbounded in this class.
Thus a method given only these saved data cannot replace a straddling
derivative interval by a certified sign using its representation label.
\end{theorem}

The capacity asymmetry matters. At $a=0$ the support supplies a lower
bound on the inward derivative, which can certify optimality when positive;
it does not supply its upper bound. At $a=A$ the opposite statement holds.
An arbitrary approximate supporting slope can have a different sign from
the actual derivative at a capacity boundary. The supplement gives an
exact rational counterexample and verifies the sharpness statements.

\subsection{A replacement which pays no second-action oracle when resolved}

Let $[g_-,g_+]$ be a valid finite version of
\eqref{eq:gradientinterval122}. For a feasible displacement
$d\in[-a,A-a]$, define
\begin{align}
 \overline m(d)&=\max\{g_-d,g_+d\}+\frac L2d^2,\label{eq:robustmodel122}\\
 \underline m(d)&=\min\{g_-d,g_+d\}+\frac\mu2d^2.\nonumber
\end{align}
The robust smooth-descent proposal is the minimizer of
$\overline m$: it is zero if $g_-\le0\le g_+$,
$-\min\{a,g_-/L\}$ if $g_->0$, and
$\min\{A-a,-g_+/L\}$ if $g_+<0$.
The executed finite menu uses the six multiples
$1,1/2,1/4,1/8,1/16,1/32$ of this proposal. Among the candidates passing all
three tests below, select the one with largest checked lower gain, with ties
resolved by the smaller realized action and then the menu index. Retain all
six attempted calculations and their costs. A smaller step can pass the
objective-width test when the unconstrained robust step cannot. This rule is
optimal only over the recorded menu under that certified-gain criterion;
it is not a global policy optimizer. The realized action must be feasible
after rounding. Recompute the
bounds at that realized action; an ideal unrounded minimizer does not
license the rounded policy.

The actual future controller is unchanged and satisfies
$J_x(b)=q(b)+\eta_x(b)$ with $0\le\eta_x(b)\le\delta$
from \eqref{eq:future121}. No regularity of $\eta_x$ is imposed.
For the realized $b=a+d$, put
\begin{align}
 U_b&=U+\overline m(d)+\delta,\nonumber\\
 L_b&=W+\min\{p_-d,p_+d\}+\frac\mu2d^2,\label{eq:candidateinterval122}\\
 B&=\min_{s\in\{p_-,p_+\}}\ 
       \min_{u\in[-a,A-a]}\{W+su+\tfrac\mu2u^2\}.
 \nonumber
\end{align}
A stronger previously verified original-optimum lower bound can replace
$B$ by its maximum with that bound.

\begin{theorem}[Original-policy gain and qualification from the same payload]
\label{thm:payloadsafeguard122}
The constructed quantities satisfy
\begin{align}
 J_x(b)-J_x(a)&\in
   [\underline m(d)-\delta,\overline m(d)+\delta],\label{eq:actualcontrast122}\\
 L_b\le J_x(b)&\le U_b,\qquad
 0\le J_x(b)-V_0(x)\le U_b-B.\label{eq:payloadqualification122}
\end{align}
Consequently a strictly negative upper endpoint in
\eqref{eq:actualcontrast122} certifies actual economic improvement.
When in addition $U_b-B\le\tau$ and $U_b-L_b\le\omega$, the replacement
satisfies the same complete-policy regret and objective-width requirements
as the incumbent. These three tests use no new original-risk value at $b$.
They do require the original paid payload, its checked smoothness bound,
all interval arithmetic, action rounding, and their computational charges.
If a test is unresolved, retain the qualified incumbent or enter the
separately budgeted original-model fallback.

A finite sequence of accepted replacements preserves feasibility and
qualification, and the sum of certified gains
$-\overline m(d)-\delta$ lower-bounds the final actual gain.
The statement is uniform in state where the premises hold; it is neither
a fitted-loss statement nor a population adoption conclusion.
\end{theorem}

For an interior anchor, an exact supplied support slope, and the exact
untruncated interval in \eqref{eq:gradientinterval122}, the robust upper
model is also sharp at each fixed displacement: the worst derivative
endpoint and its constant-curvature-$L$ witness attain it. This sharpness
does not extend automatically to a wider outward or finite-probe enclosure,
or after additional investment-specific information is imposed. Failure
of the implemented finite menu therefore records an unresolved calculation;
it does not prove that every method using the original formulas must refuse. A neural proposal can locate the first certificate evaluation, but
a spline, polynomial, matched hinge, or online method receives exactly the
same conversion and safeguard.

\subsection{Decision precision, finite refinement and priced stopping}

The precision of a supporting value and that of its action information have
different units. At an interior action with
$r=\min\{a,A-a\}>0$ and $e\le cr^2$, the derivative interval width is
at most
\begin{equation}\label{eq:precisionwidth122}
 (p_+-p_-)+2\sqrt{2(L-\mu)e}.
\end{equation}
This is the bound for the exact infimum in \eqref{eq:psi122}.
A finite list of feasible rational $h$ values remains valid but can be wider;
its actual checked radii, including any smallest-step floor, must be used.
Thus reducing a support-value error by a factor four reduces this exact
component of action uncertainty by only a factor two. Conversely, a certified
marginal tolerance $\epsilon$ requires both the slope enclosure and the
support gap to be allocated explicitly; a small critic fitting loss does
not supply either quantity.

Suppose a declared refinement ladder at this same action guarantees
$e_k\le e_0 4^{-k}$ and slope width at most $w_0 2^{-k}$, at counted
incremental work $C_k$. These are oracle assumptions to be verified for
the implemented original-risk calculations. Put
$D=w_0+2\sqrt{2(L-\mu)e_0}$.
Once the interior condition is reached, the exact-infimum derivative
width is at most $D2^{-k}$. A finite-radius implementation does not inherit
this rate without a corresponding radius-approximation guarantee. If $g=q'(a)$ and
\[
 G_L(a,g)=\max_{d\in[-a,A-a]}\{-gd-Ld^2/2\},
\]
the computed unrounded descent has guaranteed $q$ improvement at least
$G_L(a,g)-AD2^{-k}$. Hence a true smooth-model margin
$G_L(a,g)>\delta+\gamma$ resolves its gain test once
$AD2^{-k}<\gamma$, with total actual acquisition work
$\sum_{j\le k}C_j$. Qualification width is an additional test.
Even at exact derivative information, the Taylor interval can retain width
$(L-\mu)d^2/2+\delta$ at the proposed action. A ladder that cannot meet
$\omega$ from that information must acquire a candidate value or fall back;
it must not declare success from the gain test alone.

There is a narrow but exact priced stopping region. Suppose the current
qualified policy has certified original regret at most $K$, one unit of
additional computation costs $\lambda$ units of original objective, and
any allowed continuation of the solver must first pay at least $c_*>0$
new units. Exclude reuse of that information for other decisions from this
single-decision comparison. If $\lambda c_*\ge K$, stopping is weakly
optimal against every such continuation: even the exact economic optimum
cannot save more than $K$. This is a dominance rule for a stated compulsory
cost, not a general optimal-stopping theorem inferred from historical clock
timings. Actual cost intervals, possible reuse, and the complete comparative
menu remain part of an adoption study.

\section{The information in a complete Bellman support record}
\label{sec:fullsupport125}

A paid Bellman support asserts one inequality on the entire feasible action
domain. Retaining only its value and a recovered derivative interval loses
part of that assertion. With an uncertain supporting slope, the support
assertion must still hold on the whole domain. We give an exact finite
convex representation of this restriction. This advances the joint interpolation
calculation in Section~\ref{sec:joint124}, whose information class explicitly
excluded the raw supports. It does not identify a capability exclusive to
neural representations. Smooth convex interpolation has an established
literature \citep{THG2017Interpolation122}; the result below identifies the
additional constraints licensed by the original Bellman payload.
Optimal lower bounds based on collected first-order oracle information
also predate this construction \citep{FloreaNesterov2024Oracle125}.
Our contribution is the completion of the declared NBO payload on its
finite constrained action domain, a constructive sufficiency proof, and
its use in original-policy certificates. A general new theory of optimal
oracle bounds is not claimed.

Fix one state, original risk law and optimal-future objective
$q=Q_x$ on $[0,A]$, where $A>0$. Assume $q$ is continuously
differentiable, $\mu$-strongly convex, and has $L$-Lipschitz derivative,
$0<\mu\le L$. Retain the paid value and derivative intervals, and
the proved derivative cap $|q'|\le M$. For each original support record
$j$, the information also asserts that there is a single number
$p_j\in[p_j^-,p_j^+]$ such that
\begin{equation}\label{eq:globalpayload125}
 q(x)\ge S_j(x;p_j):=
 W_j+p_j(x-a_j)+\frac\mu2(x-a_j)^2
 \qquad(0\le x\le A).
\end{equation}
Its anchor, lower level and slope interval must come from the same
original calculation. A fitted derivative is not a substitute. The
quantifier in \eqref{eq:globalpayload125} is ``there exists one slope
valid at every action,'' not a new slope for each action. Exact slopes
are included by setting $p_j^-=p_j^+$. The original finite-risk support
record has this quantifier; its outward slope interval encloses the
single slope of that supporting objective.
For convex $h$ below, enforcing the entire robust pointwise lower envelope
is equivalent to existence of such a shared slope. The shared variable is
an exact convex representation of that global requirement, not an assertion
of additional information beyond the full robust obstacle.

Let $\mathcal Q_{\rm full}$ be exactly the nonempty class satisfying
these properties and all retained value and derivative intervals.
There is at least one finite value interval. Additional formulas of the
investment model are not included unless explicitly encoded. Set
$K=L-\mu$ and $h(x)=q(x)-\mu x^2/2$. A support becomes the affine
function
\begin{equation}\label{eq:affinesupport125}
 \ell_j(x;p_j)=W_j-p_ja_j+\mu a_j^2/2+(p_j-\mu a_j)x.
\end{equation}
Thus $h'$ is nondecreasing and $K$-Lipschitz, and $h\ge\ell_j$.

For $\ell\le0\le r$ define the bounded-domain conjugate
\begin{equation}\label{eq:boundedconjugate125}
 \Phi_K(v;\ell,r)=\max_{\ell\le t\le r}
              \{vt-Kt^2/2\}.
\end{equation}
When $K>0$, put $t_*={\rm proj}_{[\ell,r]}(v/K)$; then
$\Phi_K=vt_*-Kt_*^2/2$ and $\partial_v\Phi_K=t_*$.
When $K=0$, $\Phi_0=\max\{v\ell,vr\}$.
These formulas are exact rational calculations for rational inputs.

\subsection{A finite convex program that retains every support}

For trial actions $a,b$, form the sorted distinct nodes
$0=x_0<\cdots<x_n=A$ from the endpoints, every paid value and
derivative location, every support anchor, and $a,b$.
Write $f_i=h(x_i)$ and $g_i=h'(x_i)$. Transform all paid
value and derivative bounds accordingly; impose the original derivative
cap $|g_i+\mu x_i|\le M$ and finite value boxes obtained by integrating
that cap from any paid value interval. On each adjacent interval put
\[
 d_i=x_{i+1}-x_i,\qquad s_i=g_{i+1}-g_i,\qquad
 w_i=f_{i+1}-f_i-g_id_i.
\]
For $K>0$ retain the joint interpolation inequalities
\begin{equation}\label{eq:adjacentfull125}
 0\le s_i\le Kd_i,\qquad
 s_i^2\le2Kw_i,\qquad s_i^2\le2K(s_id_i-w_i).
\end{equation}
For $K=0$ use $s_i=w_i=0$. Include one shared variable $p_j$ in
its supplied interval for each support, and add, for every node and support,
\begin{equation}\label{eq:obstaclefull125}
 f_i-\ell_j(x_i;p_j)\ \ge\
 \Phi_K(p_j-\mu a_j-g_i;-x_i,A-x_i).
\end{equation}
Denote the resulting finite feasible set by $\mathcal P_{\rm full}(a,b)$.
It has a linear number of adjacent interpolation constraints and
$(n+1)m$ support inequalities for $m$ retained supports. Duplicate
observations are constraints at one node. Both domain endpoints are
present even when no measurement was bought there.

\begin{theorem}[Exact interpolation above complete Bellman supports]
\label{thm:fullinterpolation125}
A list $(f_i,g_i,p_j)$ belongs to $\mathcal P_{\rm full}(a,b)$ if
and only if there is a function in $\mathcal Q_{\rm full}$ realizing
the node values and derivatives, with these same $p_j$ satisfying
\eqref{eq:globalpayload125}. One such interpolant has a continuous
piecewise-affine derivative. The feasible set is nonempty, compact and
convex. Consequently the exact worst compatible original contrast is
\begin{equation}\label{eq:fullcontrast125}
 U_{\rm full}(a,b)=
 \max_{\mathcal P_{\rm full}(a,b)}
 \left\{f_b-f_a+\frac\mu2(b^2-a^2)\right\}
 =\max_{q\in\mathcal Q_{\rm full}}\{q(b)-q(a)\}.
\end{equation}
The extrema are attained. Removing the support inequalities enlarges
the information class, so its exact upper contrast cannot be smaller
than \eqref{eq:fullcontrast125}. The reversed optimization gives the
exact lower contrast. These are information bounds under the stated
interface; they are not oracle-complexity bounds for unrestricted access
to the original economic formulas.
\end{theorem}

The proof supplies a constructive sufficiency argument, rather than
assuming that nodewise lower bounds imply a global support. On an adjacent
interval the construction uses derivative slopes $K,0,K$ for $h$,
with their lengths fixed by the endpoint jets and the required area.
The extra term in \eqref{eq:obstaclefull125} is essential. At an
interior node a positive supporting-slope discrepancy requires a
positive value gap even when $f_i\ge\ell_j(x_i;p_j)$ alone holds.

\subsection{Verified contrasts without an exact optimizer}

Let $z$ collect $f,g,p$ in a valid finite rational box, and write all
constraints as $c_j(z)\le0$. The adjacent constraints are convex
quadratics. Each support constraint is convex because it is a convex
function of an affine argument minus an affine function. For $K>0$
it is continuously differentiable, including at its clipping points.
At $K=0$ choose any valid subgradient of its maximum of two affine
functions. Equality constraints may be represented by two inequalities.

For the minimization objective
$H=f_a-f_b-\mu(b^2-a^2)/2$, choose any rational point $z_0$,
nonnegative rational multipliers $y_j$, and valid subgradients
$v_j\in\partial c_j(z_0)$. Put
\begin{align}
 \mathcal L_0&=H(z_0)+\sum_jy_jc_j(z_0),\nonumber\\
 R&=\nabla H+\sum_jy_jv_j,\nonumber\\
 B&=\mathcal L_0+
 \sum_k\min\{R_k(\underline z_k-z_{0k}),
                         R_k(\overline z_k-z_{0k})\}.
 \label{eq:fullresidual125}
\end{align}
Then $q(b)-q(a)\le-B$ for every compatible original objective.
The formula requires neither primal feasibility nor stationarity of
$z_0$, nor strong duality. All terms are independently checkable by
rational arithmetic. A solver status alone is not a certificate.
An implementation intersects this upper bound with all earlier valid
same-objective bounds, so an unhelpful numerical dual proposal cannot
weaken its returned certificate. The arithmetic, failed proposals and
all retained supports remain charged.

For the unchanged implemented future $J=q+\eta$, $0\le\eta\le\delta$,
an upper contrast $u$ at a distinct realized action licenses a policy
change only when $u+\delta<0$. Absolute bounds on the same $q(b)$,
and a lower bound on its original optimum, must separately establish
the declared regret and objective-width requirements. At identical
actions the actual contrast is zero. No regularity of $\eta$ is
assumed, and the optimization of $q$ is not relabeled as optimization
of an arbitrary discontinuous $q+\eta$.

\paragraph{An exact decision consequence and its comparator.}
The supplement gives a rational example in which omitting a paid support
leaves the exact upper contrast for the action change $1\to0$ at
$-1/4$, whereas retaining it gives $-5/8$. The additional guaranteed
original gain is $3/8$ under the same disclosed payload; arbitrary future
error adds $\delta$ to either contrast. A cheap direct evaluation of the
support at the incumbent attains the stronger bound in this example.
It therefore demonstrates the economic loss from omitting support
information, and does not establish a computational advantage over a
conventional method that retains that information. A separate exact
counterexample shows why checking a support only at interpolation nodes
does not ensure its validity between nodes. The new theorem supplies
the finite exact program for all coupled records, while the present
original-economy experiment executes the separately specified acquisition
service. An algebraic information result is not relabeled as an executed
conic economic comparison.

\section{Paid action resolution and policy certificates}
\label{sec:paidresolution125}

The preceding theorem describes what a complete paid record permits.
It does not promise that the record resolves the decision or that the
finite program is economical to solve. The original-economy study therefore
also tests a declared service which buys new original-risk supports when
its inherited qualified policy remains unresolved. It uses the retained
rounded-midpoint contraction of Proposition~\ref{prop:midpoint123}, the
same original Bellman oracle, and the existing full-record bounds for
each actual acceptance. Midpoint localization is not introduced as a new
root-finding principle. The additions here connect its required support
precision, a complete work bound, the actor's initial economic certificate,
and the separate precision of the returned original policy.
The allocation of inexact-oracle accuracy is an established numerical
optimization problem \citep{DGN2014Inexact122,VanDesselGlineur2024}.
The work statement below is a closed-form solution of the particular
guaranteed-width planning problem specified here; its original-risk and
complete-policy premises remain part of the result.

An actor with original regret certificate
$J(a_0)-\min q\le K_0$ supplies the initial optimizer interval
\begin{equation}\label{eq:warmbracket125}
 I_0=[a_0-\sqrt{2K_0/\mu},a_0+\sqrt{2K_0/\mu}]\cap[0,A].
\end{equation}
An outward radius and intersection with previously certified derivative
brackets preserve validity. Every conventional actor with the same
certificate receives the same reduction. A small training loss without
that original-regret certificate does not shrink the interval.
The executed warm service uses its retained derivative-envelope bracket;
\eqref{eq:warmbracket125} is an additional valid initialization rule,
not an unrecorded change to that frozen implementation.

\begin{theorem}[A complete precision and work account for paid localization]
\label{thm:paidwork125}
Fix a relative rounding allowance $0\le\xi<1/2$ and
$\gamma=1/2+\xi<1$. Let $w_0$ bound the width of $I_0$, and choose
planned widths $w_{k+1}\in[\gamma w_k,w_k]$. Query an actually
feasible rounded midpoint of the current certified interval with rounding
error at most $\xi$ times its actual width. Suppose an oracle can return
an original support with level gap $e=U-W$ and slope width $v=p^+-p^-$
satisfying, for the requested next width $t=w_{k+1}$,
\begin{equation}\label{eq:jointprecision125}
 h=\chi t,\qquad
 \chi=\begin{cases}\min\{1/2-\xi,\mu/(4K)\},&K>0,\\
                         1/2-\xi,&K=0,\end{cases}
 \qquad e\le\frac{\mu\chi t^2}{4},\qquad v\le\frac{\mu t}{4}.
\end{equation}
Use a query only while the actual interval width exceeds $t$;
otherwise its planned contraction is already achieved. Then both
displacements $\pm h$ are feasible from the midpoint, and the finite
support conversion gives a derivative interval of width at most $\mu t$.
The new optimizer interval has width at most $t$.

Assume, as a separately verified oracle premise, that all computations
needed for one such support, including its unsuccessful precision levels,
cost at most $C_0+C_1t^{-p}$, with $C_0,C_1\ge0$ and $p>0$.
For a target $0<w_*<w_0$, put
\[
 N=\left\lceil\frac{\log(w_0/w_*)}{\log(1/\gamma)}\right\rceil,
 \qquad w_j=w_*\gamma^{-(N-j)},\quad 1\le j\le N.
\]
After at most $N$ completed precision-qualified probes, the optimizer
interval has width at most $w_*$. Their total work is bounded by
\begin{equation}\label{eq:geometricwork125}
 NC_0+C_1w_*^{-p}\frac{1-\gamma^{pN}}{1-\gamma^p}.
\end{equation}
Among all nonincreasing planned-width schedules with guaranteed
contraction floor $\gamma$, final width at most $w_*$ and this same
declared nonincreasing per-probe cost bound, this schedule minimizes the
sum of those bounds. It is an optimum for the stated certificate-planning
problem, not for all possible original-model algorithms or actual
shared-host elapsed times. Querying every probe at the terminal target
instead gives the bound $NC_0+NC_1w_*^{-p}$.

To return a policy, let a paid original-value interval at a feasible
$b$ be $[W_b,U_b]$, of width $e_b$, and let
$r_b=\max\{|b-\underline a|,|b-\overline a|\}$ for the final
optimizer interval. Then
\begin{align}
 |b-a_q^*|&\le r_b,\nonumber\\
 W_b\le J(b)&\le U_b+\delta,\nonumber\\
 0\le J(b)-\min q&\le\delta+Mr_b.
 \label{eq:terminalpolicy125}
\end{align}
If the optimizer is certified to be stationary, $Mr_b$ can additionally
be replaced by $Lr_b^2/2$. This regret bound does not require the value
width $e_b$: a certified original optimizer interval already controls
the economic action error. In particular, an interval strictly inside
$(0,A)$ certifies stationarity. A returned action must separately satisfy
$e_b+\delta\le\omega$ and the declared regret tolerance. Any fresh
value calculation needed for this final gate is charged in addition to
\eqref{eq:geometricwork125}, as are actor acquisition, the inherited
service and the certificate algebra. The theorem does not turn a narrow
optimizer interval into an unpaid absolute value certificate.
\end{theorem}

The proof uses a finite feasible radius in
\eqref{eq:jointprecision125}; it does not assume that an implemented
finite-radius conversion attains a real-valued infimum. The oracle-work
premise includes the required slope enclosure and original continuation
accuracy. A terminal-solver floor, a maximum risk resolution, an action
lattice floor or a resource cap can prevent it from holding. The executed
study records such cases as unresolved and retains the last qualified
policy. It neither fits a favorable power law to the realized times nor
claims \eqref{eq:geometricwork125} as an empirical complexity theorem
without verifying that premise.
The experiment executes its frozen finite integration ladder and records
achieved contractions. It does not claim to execute the ideal
cost-minimizing planned-width schedule in the theorem.

Every accepted replacement still requires its own strictly negative
upper bound on $J(b)-J(a_0)$ under the same future controller, as well
as the original regret and width gates. The study targets a stricter
regret tolerance as an additional outcome; a failure to attain that
target does not revoke the already qualified incumbent. Failed and
rejected queries remain in its complete service bill. These distinctions
connect information acquisition to an economic action rather than to an
internal fitting error.


\section{An independent original-economy decision-service study}
\label{sec:experiment125}

\subsection{Economic object, comparisons and prospective identity}

The new computation uses exactly the two-dimensional investment primitives
in Section~\ref{sec:continuous95}. Both risk levels $\alpha=1/2,3/4$
are included. The implemented continuation controller is unchanged, with
$q_x\le J_x\le q_x+\delta$ and $\delta=15/2^{24}$. A service qualifies
only when its outward original-regret bound is at most $1/256$ and its
absolute policy-value interval has width at most $1/8192$. A separately
declared, stronger regret target is $1/65536$. Passing the ordinary service
gate does not count as meeting that stronger target.

The protocol and complete executable source closure were committed at
\nolinkurl{a94e1f43bf64d8cb90af21e4b53d07335c615576} and all 338
published file entries were read back before any new economic evaluation.
The four states, in their fixed order, are
\[
 (1/4,1/4),\quad(1/4,3/4),\quad(3/4,1/4),\quad(3/4,3/4).
\]
This is the full deterministic census of these states. It does not estimate
an expectation over a state distribution and has no sampling confidence
interval. The 56 actor files retain their original acquisition commit
\nolinkurl{796c5f2badf68b051ac6c144884c0e61b87faf5e}. The two state-mesh sizes, five and nine points per coordinate,
give 25 and 81 acquisition states. Thirteen original-objective contexts
per state give 325 and 1,053 contexts per risk. Each acquired design
is shared across seven representations and value versus marginal fitting. The ReLU and ordinary hinge actors use the
same features and information; the hinge account includes its required
parent acquisition. No actor is refitted or selected after these outcomes.

Every actor and each risk-specific online control uses four complete
services. Each service independently pays for the same strong fused
original-model prefix, including qualification, fallback, retained
supports and rejected finite-menu proposals. The control returns that
completed service. The warm adaptive service initializes an optimizer
bracket from all its paid derivative envelopes; the cold adaptive service
begins from the entire feasible interval. The fixed-precision warm service
uses 512 risk bins at each new center. Two further conventional controls
evaluate a 33-action uniform dynamic-programming grid under the original
risk law. The full catalogue has 234 cohorts and 936 primary decisions.

\subsection{Action acquisition, stopping and independent checks}

For adaptive precision the fixed risk-bin ladder is
$32,64,128,256,512,1024,2048$. The service attempts at most eight
new action centers and at most 32,768 additional terminal midpoint or
endpoint contexts. Each $n$-bin query is charged $2n+1$ such contexts,
as well as its actual terminal iterations, risk operations and elapsed
time. Coarse queries discarded after refinement remain charged. The
completed prefix is additional to this cap. Actual midpoint actions lie on
the inherited feasible 30-bit lattice; their observed rounding displacement
enters the contraction inequality.

An action replaces the completed prefix only when a strictly negative
upper bound on its original-policy contrast, its original-regret bound and
its absolute-width bound all pass. The future-policy allowance $\delta$
is charged before acceptance. Strict gains that miss the stronger target
remain target failures. A maximum precision, exhausted budget, unresolved
derivative or failed numerical calculation retains the last safe action and
its full cost. The grid control evaluates all 33 actions at 128 bins,
refines its predeclared best-upper-value action when needed, and pays a
complete online fallback if no grid action qualifies.

The first verification replays the saved rational certificate algebra,
feasibility, brackets, budgets, accepted gains and exact matched-feature
identities. A separate implementation reconstructs the continuous-law
economics from the retained rational primitives, without importing the
production economic evaluator. It independently encloses the optimized-future
objective $q$ and uses the already proved same-controller allowance $\delta$
to enclose $J$; it does not replay every rounded terminal-controller action.
Every distinct deployed and prefix action
receives its own eight-bin evaluation. Unresolved checks are refined at
32 and 128 bins in a fixed hash order under a total 65,536-context cap,
with 32 terminal bisections per context. Interval containment, confirmation
of an accepted gain, overlap and contradiction are reported separately.
An unresolved check is not a successful independent confirmation. The
terminal witnesses and full actual bisection ledger are retained.

The exact mathematical fixtures check the complete-support witness and
the necessity of between-node obstacle constraints. An independent
breakpoint calculation also checks the interpolation construction using
rational arithmetic. These calculations validate finite algebraic examples;
they are distinct from the original-economy decisions and from a formal
proof-assistant verification.

\subsection{Complete costs and the adoption question}

Exactly three separate serial timing processes execute all 234 services at
the first declared state, with a fixed rotation of cohort order. Their 702
decisions must reproduce the primary actions, certificates, operation
counts and stopping reasons. The host is shared and its frequency and load
are uncontrolled. The observed clock range is descriptive. Each measured
service includes model loading, all prefix and new work, output and
bookkeeping, plus its allocation of the measured parent-process residual.
Historical acquisition and fitting allocations remain identified as
historical measurements. Missing interrupted historical process clocks
remain unknown. The new mathematical checks, reproduction and independent
economic evaluation enter the research-work ledger separately.

The priced adoption comparison is restricted to this first-state timing
workload. For each risk and each $N\in\{1,32,1000,10^6\}$, it compares
\[
        J+\lambda\{t+K/N\},\qquad \lambda\ge0,
\]
where $t$ is complete per-decision time and $K$ is the retained acquisition
bill. Every traditional service is included. Intersecting the linear
inequalities over all comparators gives sufficient adoption regions from
own upper and rival lower cost bounds, and optimistic outer regions from
the reversed bounds. The price $\lambda$ is a declared exchange rate
between objective units and seconds, not an empirically calibrated value
of computation. Four-state economic results are not combined with
first-state timing as though they measured the same workload.

This experiment tests scalar-action localization and continuous-risk
precision in the original two-date economy. It does not establish scaling
in the number of controls, independent innovations, state coordinates or
dates, and it does not recalibrate the economic primitives. The complete
source, original protocol, all attempts and hashes allow those restrictions
to be assessed directly. Protected R123 and R124B experiments retain their
own source and execution identities; this study neither completes them nor
substitutes its results for their outcomes.

\section{Original-economy action acquisition}
\label{sec:results125}
The prospectively frozen R125 census evaluates the four dyadic states $(1/4,1/4)$, $(1/4,3/4)$, $(3/4,1/4)$ and $(3/4,3/4)$ under both original risk levels. All 56 inherited actors, both original-model online controls, and both conventional action-grid controls are retained. These are fixed cases, not a population sample. Every paired economic comparison holds the risk parameter and original objective fixed; the equally weighted table summaries across risk/state cases are descriptive, not a representative household's welfare or a population expectation. The 936 original-service decisions all satisfy regret $1/256$ and value width $1/8192$; the separately specified stronger regret target is $1/65536$. The future controller and continuous-risk economic objective are unchanged.
\begin{table}[t]
\centering\small
\caption{Prospective original-action acquisition census}
\label{tab:acquisition125}
\begin{tabular}{lrrrrr}
\hline
Service & Decisions & Target & Gains & New contexts & Total seconds \\
\hline
Strong fused control & 232 & 147 & 0 & 0 & 6.156 \\
Warm adaptive & 232 & 231 & 84 & 37,557 & 15.031 \\
Cold adaptive & 232 & 229 & 83 & 73,015 & 31.977 \\
Warm fixed 512 & 232 & 231 & 84 & 181,425 & 13.376 \\
Uniform 33-action DP & 8 & 6 & 0 & 83,216 & 3.556 \\
\hline
\end{tabular}
\par\smallskip\parbox{0.96\textwidth}{\footnotesize Strict gains are relative to the completed strong fused prefix. New contexts count every paid midpoint and endpoint at every attempted precision; all prefix and rejected-candidate work is additionally charged. Total seconds sum complete four-state cohort clocks and their equal share of the observed parent residual. This is one execution, not a timing confidence interval; historical acquisition is separately charged in the complete catalog. Uniform DP does not claim a strict gain over a prefix.}
\end{table}
Warm adaptive and fixed-512 localization both attain the stronger target in 231 of 232 services and certify 84 additional strict original-policy gains. Adaptive precision uses 79.30 percent fewer extra terminal contexts, but requires 12.38 percent more complete elapsed time in this primary execution. Its extra interval algebra and coarse queries therefore matter economically. The two services return identical action bits in 205/232 pairs and identical cost intervals in 148/232; all 232 target and acceptance indicators agree. Thus equal target attainment does not imply identical policies. One singleton-bracket case does not pass the stronger certificate gate for either method; cold localization retains three maximum-center stops. No exceptions were recorded in the enhancement services.
The saved clocks locate the computational tradeoff. Query and gradient-recovery time is 3.195700 seconds for warm adaptive and 5.348735 seconds for fixed 512, whereas the remaining time within the service is 5.469696 and 1.563042 seconds, respectively. This remainder includes setup, support-case and candidate calculations, selection and bookkeeping; it is not an isolated algebra timer.
\begin{table}[t]
\centering\small
\caption{Original economic objectives and decision-loss bounds}
\label{tab:economics125}
\begin{tabular}{lrlrr}
\hline
Service & $n$ & Mean original $J$ interval & Mean regret & Max regret \\
\hline
Strong fused control & 232 & $[0.5322899,0.5322980]$ & 0.00046336 & 0.00364741 \\
Warm adaptive & 232 & $[0.5320457,0.5320480]$ & 0.00000372 & 0.00001781 \\
Cold adaptive & 232 & $[0.5320473,0.5320497]$ & 0.00000609 & 0.00015652 \\
Warm fixed 512 & 232 & $[0.5320461,0.5320473]$ & 0.00000271 & 0.00001781 \\
Direct original online & 8 & $[0.5320452,0.5320464]$ & 0.00000111 & 0.00000112 \\
Uniform 33-action DP & 8 & $[0.5320498,0.5320509]$ & 0.00000977 & 0.00003496 \\
\hline\end{tabular}
\par\smallskip\parbox{0.96\textwidth}{\footnotesize The first four services use the same 232 actor/risk/state cases. The direct online and uniform-DP controls use the same four-state and two-risk weights, without averaging over fitted representations. Lower value endpoints are rounded downward; upper endpoints and displayed regret bounds are rounded upward using exact rational arithmetic. Regret is bounded relative to the original optimized-future objective; these are finite-census bounds, not confidence intervals.}
\end{table}
The singleton case isolates a remaining implementation issue. It is the risk-$1/2$, budget-5, marginal-fitted cubic actor at $(1/4,3/4)$, returning $a=3/16$ with the exact original-optimizer bracket $[3/16,3/16]$. Its retained level-based certificate is $1.7807343\times10^{-5}$, above the stronger target. The terminal action-radius bound of Theorem~\ref{thm:paidwork125} would give $\delta$ at this already optimal original action. That terminal gate is not implemented in the frozen service; we therefore retain the failed target and the count 231/232. The diagnostic distinguishes action uncertainty from absolute-value precision without changing an observed outcome.
All 702 decisions in three serial timing passes reproduce the primary action, certificate, stopping reason and operation counts. The saved-algebra audit executes 16,162 assertions, including 128 matched ReLU/ordinary-hinge pairs. The independent rational original-law audit confirms 850 of 936 production intervals by containment and 193 of 251 accepted gain bounds, with 0 contradictions. Any unresolved check remains explicitly unresolved. The audit uses 52,195 terminal contexts and retains all original continuous-risk primal--dual witnesses.
\paragraph{Separately frozen post-primary independent verification.}
After the original capped audit, a separately published protocol selected both actual final and completed-prefix actions for every unresolved check. Its complete union contains 110 distinct state/risk/action keys; each received one 512-bin independent original-law evaluation. This is an explicitly post-primary verification follow-up, with no changed economic state, model, policy, training sample or production target.
Combining the additional independently valid intervals with the original audit confirms 934/936 production interval containments and 251/251 accepted gain bounds, leaving 2 and 0 unresolved, respectively. There are 0 contradictions. The original capped result remains 850/936 and 193/251, separately reported. No frozen success count is relabeled.
The additional audit requests 112,750 terminal contexts, reuses 28,270 exact original witnesses at identical rational states and risks, and computes 84,480 new terminal contexts with 2,586,336 actual bisection iterations. Its complete parent-process clock is 1459.236044 seconds, charged in addition to the original audit. These are research-verification costs; neither production-service clocks nor the first-state adoption comparisons are re-timed or revised.

Complete-cost neural adoption is assessed only on the first-state workload for which three complete timing processes were observed. Of 128 prespecified risk/budget/loss/service and amortization comparisons, 128 sufficient and 0 optimistic price regions are empty. These conditional timing comparisons are not statistical confidence statements. Full four-state original-value intervals, regret and action-distance bounds, acquisition charges, complete measured service costs and failures are provided in the development companion and machine-readable catalog.
All 58 learned/online prefixes already meet the stronger target at that first state, so the three repeated timing passes execute zero additional localization queries in each of the four prefix-based services. They therefore do not measure the precision-acquisition tradeoff observed in the other census states. The complete primary clocks above retain that work, with one observed execution. The nonempty optimistic regions are the original interval-only outer sets; the policy-identity analysis below supplies additional equality information without revising those original classifications.
The support-complete interpolation theorem and its exact rational witnesses have a separate evidential role: the production services above use original-support and derivative-envelope acquisition, not an implemented full-support conic solver. The strict support witness verifies an economic contrast improvement relative to omitting paid information; its endpoint bound is also attained by cheap direct support evaluation. Neither result supplies a neural-specific separation from matched hinge features.

\section{Economic interpretation and the next decision problem}
\label{sec:discussion125}

The economically relevant output is a feasible investment and its value
under the stipulated recursive preference. A neural continuation can be
useful in this account in two ways: it can propose an action that reduces
original-policy loss, or it can reduce the complete information cost needed
to reach a specified economic tolerance. A reduction in residual span,
critic training error or an optimizer's internal objective establishes
neither consequence without the accompanying action and cost evidence.

Table~\ref{tab:economics125} puts the improvement in original economic
units. Subtracting its outward-rounded intervals, the warm adaptive
service reduces mean cost relative to the completed fused control by
between $0.0002419$ and $0.0002523$ on this census. Its mean cost minus
that of the direct original online control lies between $-0.0000007$
and $0.0000028$. The comparison therefore does not certify an economic
improvement over that strong conventional control. These are objective
units in the specified model, not an estimated monetary benefit.
The eight direct-online cases consume $0.3750481$ complete primary
seconds in total, or about $0.0468810$ seconds per decision, with no
training acquisition bill and all eight stronger targets attained.
They are a subset of the fused-control row, not additional observations
or additional work. This one observed clock is descriptive; the full
decision cost and adoption conditions remain part of the comparison.

Complete Bellman supports make this distinction mathematical. A paid
calculation asserts a restriction on all feasible actions. Its information
should be reused in its full form when comparing decisions. The exact
finite representation shows when a numerical certificate is the strongest
one compatible with that declared record, and when apparently adequate
nodewise checks admit an impossible global support. Its rational action
example strengthens a guaranteed gain by $3/8$ in the specified
information class. A conventional direct-support calculation also obtains
that example's gain. The theorem establishes the full compatibility
problem; practical superiority of a complete-support solver over strong
conventional bounds requires a separate executed comparison.

The precision account treats new information as an economic input. A
smaller optimizer bracket may warrant a tighter risk calculation, whereas
a fixed terminal precision can leave a service unable to resolve the
action. Width, regret and improvement have different gates: a narrow
derivative bracket alone does not provide an absolute value enclosure,
and an optimal-future improvement must still pay for the actual future
controller. The implemented finite ladder records these distinctions.
It preserves qualified decisions when a new calculation cannot certify
a better one; unsuccessful acquisition remains part of the cost.


A frozen service outcome exposes this distinction directly. One warm
and fixed-precision cell reaches the capacity action $A=3/16$ and a
singleton original optimizer bracket, but its stored value-derived regret
bound is approximately $1.780734\times10^{-5}$, above the declared
$1/65536$ target. A separately recorded post-freeze rational diagnostic
reconstructs the original boundary derivative from the saved supports and
proves $q'(A)<0$. Strong convexity therefore identifies $A$ as the unique
constrained optimizer. The unchanged future controller then gives regret
at most $\delta=15/2^{24}$, approximately $8.94\times10^{-7}$, by
Theorem~\ref{thm:paidwork125}. The action and its absolute value interval
do not change. This diagnostic isolates an unnecessarily weak decision
gate and supplies a precise prospective implementation correction; it does
not alter the frozen target counts, assert a new policy gain, or claim
that the improved gate was deployed in the experiment.

The representation comparison is equally explicit. An affine--ReLU circuit
and the identical native hinge function are alternative encodings of the
same acquired decision rule. Their shared geometry does not identify a
neural-specific statistical or computational effect. A convincing neural
advantage must survive comparisons with conventional representations
having the same acquired labels, features, feasible policies and work
budget, and with suitable analytic or semianalytic original-model solvers.
The retained quadratic wins and empty adoption regions are evidence about
that requirement. They are not removed when a new certificate improves
an internal bound.

The current application is a stylized constrained investment economy.
Its continuous innovation and nested tail preference make qualification
nontrivial, but it has one control, one innovation factor and two decision
dates. The derivative constants, exact interpolation and execution results
have this stated scope. Neither a black-box continuous-risk lower bound
nor a scalar support-information minimax result excludes a faster method
that uses additional primitive formulas. Nor does a conditional precision
cost law prove a uniform lower bound over all economic algorithms.

The research programme therefore has concrete remaining tests. A neural
method must acquire action-relevant continuation information more cheaply
than equally informed strong conventional methods, including an independent
refit account when its training rule changes. A multidimensional extension
must replace scalar ordering and control the cost of continuous risk,
constraints and boundary interactions. Economic adoption needs either a
calibrated decision environment or an explicit range of values of the
computational resource. A new actor's population claim requires its own
independent sample and inference, rather than reuse of an earlier actor's
confidence interval. These are extensions of the original NBO objective
and research ambition, with the present theorem, code and preserved
negative evidence specifying what a successful extension must improve.

\bibliographystyle{ecta-fullname}
\begingroup
\let\NBOBibliographyStart\thebibliography
\renewcommand{\thebibliography}[1]{\NBOBibliographyStart{#1}\raggedright}
\bibliography{references99,references110-additions,references120,references121,references122,references124,references125}
\endgroup
\end{document}
